\documentclass[10pt,a4paper]{amsart}
\usepackage[margin=19mm]{geometry}
\usepackage{amsmath,amssymb,mathtools}
\usepackage[scr=rsfs,cal=euler]{mathalfa}
\usepackage{xfrac}
\usepackage[unicode,hidelinks,bookmarks=false]{hyperref}
\newcommand{\ie}{i.e.\ }
\newcommand{\cf}{cf.\ }
\newcommand{\resp}{respectively\ }
\newcommand{\Subsection}{\subsection}
\providecommand{\nobreakdash}{\nobreak\hbox{-}}
\setlength{\emergencystretch}{2em}
\multlinegap0pt
\usepackage{amssymb}
\usepackage{csquotes}
\usepackage{bm}
\usepackage[shortlabels]{enumitem}
\newtheorem{proposition}{Proposition}
\def\WW{\mathbb{W}}
\def\calW{\mathcal{W}}
\title{The Shortest-Path distance on graphons}
\author{C\'edric Simal, Julien Petit, Timoteo Carletti}
\date{Source: arXiv:2506.14353v2 (CC BY 4.0)}
\begin{document}
\maketitle

\noindent\emph{Source and licence.} The Shortest-Path distance on graphons. Version arXiv:2506.14353v2 (CC BY 4.0), distributed by arXiv under the Creative Commons Attribution 4.0 International licence, http://creativecommons.org/licenses/by/4.0/, as recorded in the arXiv licence metadata for this exact version and shown on the abstract page https://arxiv.org/abs/2506.14353v2. Full licence notice: \texttt{licenses/arxiv-2506.14353-CC-BY-4.0.txt}.\par\smallskip

\noindent\textit{Editorial scope.} Counted unit: the article's proposition prop:isomorphic-graphon-unitary (source0.tex lines 196--198). The article's complete original proof of it is delivered with it (source0.tex lines 199--213, one \texttt{proof} environment of 15 lines). No mathematics is written, repaired or re-derived: the statement and the proof are the article's own lines, in the article's own notation, and no retained line is substituted or re-flowed.  No further source-local context is needed: every cross-reference of every retained line resolves inside the two delivered spans.  Nothing was omitted from any retained span, so the package carries no omission record.  No retained line contains a citation, so the wrapper carries no bibliography and no bibliography entry.  No retained line contains a figure environment, an image-inclusion command or drawing code, so no image is delivered and the package declares no asset dependency.  Editorial formatting only, in addition: an AMS \texttt{amsart} preamble that restates the article's own declarations and macros used by the retained lines, namely the environments \texttt{proposition} on separate counters, together with the standard packages those lines need.  The retained environments print in this document as Proposition 1 in order of appearance, whereas the article prints the same items with the numbering of its own full document.  The complete native source of the article as distributed in the arXiv e-print for this exact version is bundled unchanged as \texttt{sources/180006/source0.tex}.
\begin{proposition} \label{prop:isomorphic-graphon-unitary}
    If two graphons are isomorphic, then their adjacency operators are unitarily equivalent.
\end{proposition}

\begin{proof}
    Let $W \in \WW$ and $\varphi: [0,1]\rightarrow [0,1]$ a measure preserving bijection. We construct a unitary operator $U^\varphi$ on $L^2 [0,1]$ induced by $\varphi$ as
    $$ (U^\varphi f)(x) = f(\varphi(x)). $$
    The operator $U^\varphi$ is easily proved to be linear, and letting $(U^\varphi)^*$ denote the adjoint of $U^\varphi$, we have for all $f,g\in L^2 [0,1]$,
    \begin{align*}
        \langle (U^\varphi)^* f, g\rangle = \langle f, U^\varphi g\rangle &= \int_0^1 f(x) g(\varphi(x)) dx = \int_0^1 f(\varphi^{-1}(x)) g(x) dx,
    \end{align*}
    so $((U^\varphi)^* f)(x) = f(\varphi^{-1}(x))$, and it is clear that the $(U^\varphi)^* = (U^\varphi)^{-1} = U^{\varphi^{-1}}$, so $U^\varphi$ is unitary.

    Now, for $x\in [0,1]$, we compute
    \begin{align*}
        (U^\varphi \calW U^{\varphi^{-1}} f)(x) & = \int_0^1 W(\varphi(x), y) f(\varphi^{-1}(y)) dy = \int_0^1 W(\varphi(x), \varphi(y)) f(y) dy,
    \end{align*}
    so $U^\varphi \calW U^{\varphi^{-1}}$ is equal to the adjacency operator of $W^\varphi$, which is isomorphic to $W$. Since any graphon isomorphic to $W$ can be expressed in this form, we conclude the proof.
\end{proof}

\end{document}
